\documentclass[12pt,letterpaper]{article}

% Librerías a utilizar
\usepackage[utf8]{inputenc}	% Codificación 
\usepackage[spanish]{babel}	% Idioma
\usepackage{natbib}			% Bibliografía
\usepackage{graphicx}		% Imagenes
\usepackage{indentfirst}		% Sangría
\usepackage{amsmath, amsfonts, amssymb}	% Figuras matemáticas
\usepackage{url}							% URL
\usepackage{subfigure}
\usepackage{svg}
\usepackage{algpseudocode,algorithm}
\floatname{algorithm}{Algoritmo}
\usepackage[toc,titletoc,title]{appendix}
\renewcommand{\appendixtocname}{Anexos}
\renewcommand{\appendixpagename}{Anexos}
\usepackage{tabularx}
\usepackage{pdflscape}
\usepackage{xltabular}
\usepackage{multirow}
\usepackage{booktabs}
\usepackage[table]{xcolor}
\usepackage{colortbl}
\usepackage{array}\setlength{\extrarowheight}{1pt}
\usepackage{rotating}
\usepackage{stfloats}
\usepackage{placeins}
\usepackage{float}
\usepackage{caption}
\usepackage{enumitem}

\usepackage{fancyhdr}
\fancypagestyle{cartastyle}{
	\fancyhf{}
	\renewcommand{\headrulewidth}{0pt}
	\renewcommand{\footrulewidth}{0.4pt}
	\fancyfoot[C]{\footnotesize \color{gray} * Nota: Se utilizó el modelo Gemini 3.6 Flash (Medium) para la redacción del texto y no para el proceso de pensamiento/razonamiento.}
}

% Enlaces cliqueables en el PDF (citas, índice, figuras, etc.)
\usepackage[hidelinks]{hyperref}
\newcommand{\appendixnumberline}[1]{\makebox[5.5em][l]{#1}}



%\usepackage[nonumberlist]{glossaries}
%\makeglossaries


\tolerance=1
\emergencystretch=\maxdimen
\hyphenpenalty=10000
\hbadness=10000

\renewcommand{\theenumi}{\alph{enumi}} %Enumerar Letras minúsculas 

\usepackage[left=3cm,right=2cm,top=2cm,bottom=3cm]{geometry}

\setlength{\parindent}{2cm}	% Sangría en los párrafos
\renewcommand{\baselinestretch}{1.5}	% Interlineado

% Renombrar ciertos títulos del texto
\renewcommand\spanishcontentsname{Tabla de contenidos}
\renewcommand\spanishlisttablename{Índice de tablas}
\renewcommand\spanishrefname{Referencias}
\renewcommand\spanishtablename{Tabla}

% Inicio del documento
\begin{document}
%

% Carta de Respuesta a la Comisión Evaluadora
\input{carta_respuesta}

%%%%%%%% PORTADA %%%%%%%%%

\newpage
\vspace*{-.5cm}
% Logo institucional
\begin{picture}(18,4)(0,30)
	\put(350,-20){\includegraphics[scale=0.25]{Images/LogoUsach.pdf}}
\end{picture}

\sloppy
\thispagestyle{empty}
\vspace*{-1.6cm}

% Datos institucionales
\begin{center}
	{\bf \mbox{\large UNIVERSIDAD DE SANTIAGO DE CHILE}}\\
	{\bf \mbox{FACULTAD DE INGENIERÍA}}\\
	{\bf \mbox{DEPARTAMENTO DE INGENIERÍA INFORMÁTICA}}\\
\end{center}

\vspace{5cm}
%Título del trabajo
\begin{center}
	\Large
	\textbf{Método computacional basado en procesamiento de lenguaje natural para analizar la alineación curricular de Ingeniería Informática en Chile con estándares ACM/IEEE y necesidades de la industria}
\end{center}



% Datos personales
\vspace*{1.25cm}
\begin{flushright}
	\begin{tabular}[t]{l l}
		Estudiante:
		               & Fabián Ibarra Norambuena \\
		Profesor guía: & Manuel Villalobos Cid    \\
	\end{tabular}
\end{flushright}
\begin{center}
	\vspace{1.85cm}
	Santiago - Chile\\
	2026
\end{center}



\newpage
\pagenumbering{gobble}


\section*{Resumen}

El presente trabajo corresponde a un proyecto de \textbf{desarrollo} enfocado en caracterizar y analizar cuantitativamente las diferencias existentes en los programas de Ingeniería Informática en Chile mediante un enfoque de análisis computacional de mallas curriculares y perfiles de egreso declarados. Se propone un método computacional basado en procesamiento de lenguaje natural (Sentence-BERT) para la representación vectorial semántica de las asignaturas de las mallas y las competencias declaradas en los perfiles de egreso, junto con una estrategia tripartita de aprendizaje no supervisado (\textit{clustering} K-Means, Jerárquico y DBSCAN) para identificar patrones estructurales de distribución disciplinar y perfiles formativos entre carreras. Adicionalmente, se incorporan métricas de similitud de coseno y tasa de cobertura proyectando en un mismo espacio vectorial latente ($\mathbb{R}^d$) los dos objetos curriculares declarados, las \textit{Knowledge Areas} de los estándares internacionales de referencia (ACM/IEEE CS2013 y CC2020) y las competencias demandadas por el mercado laboral local extraídas de 200 vacantes de LinkedIn estratificadas entre empresas del sector tecnológico y empresas no-TI. Para evaluar la validez del agrupamiento y la alineación, se establece un marco de evaluación híbrido basado en métricas intrínsecas (Silhouette y Davies-Bouldin), validación externa contra las 18 áreas del estándar ACM mediante el curado manual de 747 asignaturas como etiqueta de referencia (\textit{ground truth}) calculando ARI, NMI y Pureza, análisis de estabilidad mediante bootstrap y evaluación cualitativa por un panel independiente de 3 expertos. Finalmente, el pipeline algorítmico y sus resultados analíticos se operacionalizan en un tablero de control interactivo (\textit{dashboard}) en React concebido como una herramienta de apoyo a la toma de decisiones para comités y gestores académicos.\\

\noindent
\textbf{Palabras claves}: Educación en Ingeniería, Algoritmos de agrupamiento, Procesamiento de lenguaje natural, Mallas curriculares, Minería de textos, Alineación curricular.

\newpage
\tableofcontents

% \listoffigures
% 
% \newpage
% \listoftables
% 
% \newpage

\renewcommand{\thepage}{\arabic{page}}
\setcounter{page}{1}

% Capítulos

% Entrega 1
\section{Introducción}
\label{cap:introduction}

\subsection{Contexto y Motivación}
\label{sec:context_motivation}

En las últimas décadas, el sistema de educación superior en Chile ha alcanzado una alta cobertura y diversidad institucional en la formación en tecnología \citep{subsecretaria_de_educacion_superior_evidencias_2025}. Esta expansión ha generado una marcada heterogeneidad nominal y de estructuración modular entre los planes de estudio dictados por las distintas universidades, coexistiendo programas tradicionales de ciclo largo (e.g., \textit{Ingeniería Civil en Informática}) con carreras de corte tecnológico y de ejecución (e.g., \textit{Ingeniería de Ejecución en Computación e Informática}) \citep{subsecretaria_de_educacion_superior_evidencias_2025,latuz_pina_2024}. Esta diversidad curricular responde de manera legítima a la autonomía de las instituciones de educación superior, a sus distintos sellos formativos y a sus modelos pedagógicos orientados a diversos nichos del mercado laboral.

No obstante, esta variabilidad estructural dificulta la comparación cuantitativa de la oferta nacional frente a marcos de referencia internacionales. Para responder a este desafío, organizaciones como ACM e IEEE han promulgado directrices curriculares de validez global, tales como CS2013 y CC2020 \citep{cc2020_task_force_computing_2020}. Si bien estos estándares proporcionan taxonomías claras de áreas del conocimiento (\textit{Knowledge Areas}), su adopción e integración en las mallas de las universidades chilenas ha sido desigual \citep{diaz-quezada_rediseno_2019}. A su vez, la rápida evolución tecnológica exige contrastar de forma continua las asignaturas de las mallas con la demanda técnica real del mercado laboral local.

En este escenario, los diagnósticos y rediseños curriculares han dependido tradicionalmente de revisiones manuales por comisiones académicas o encuestas de percepción cualitativas \citep{contreras_diseno_nodate,diaz-quezada_rediseno_2019}, procesos laboriosos y susceptibles a sesgos informativos. Para apoyar esta labor, se requiere un \textbf{método computacional analítico y bidireccional} que permita procesar objetivamente los datos textuales de la oferta educativa y contrastarlos dinámicamente. El carácter \textbf{bidireccional} de este enfoque abarca dos vertientes analíticas complementarias:
\begin{enumerate}[label=\alph*)]
	\item \textbf{Vertiente Académico-Normativa (Top-Down):} Evalúa cuantitativamente la densidad y cobertura disciplinar de las asignaturas de las mallas curriculares frente a las directrices de los estándares internacionales ACM/IEEE CS2013 y CC2020.
	\item \textbf{Vertiente Mercado-Laboral (Bottom-Up):} Mide la alineación semántica entre las mallas curriculares y las demandas tecnológicas extraídas en tiempo real desde 200 vacantes laborales de LinkedIn, estratificadas formalmente entre empresas del sector tecnológico (Mundo TI) e industrias usuarias no tecnológicas (banca, minería, retail, servicios públicos).
\end{enumerate}

Para asegurar un análisis delimitado y transparente, el método computacional procesa \textbf{única y exclusivamente dos objetos curriculares declarados}: (1) las \textbf{mallas curriculares} (títulos, descriptores modulares, distribución de créditos y prerrequisitos de 747 asignaturas), y (2) los \textbf{perfiles de egreso declarados} por cada una de las 16 carreras universitarias que componen el corpus de estudio (Anexo C).

\subsection{Propósito de la solución}
\label{sec:purpose}

El propósito de este trabajo es proyectar y desarrollar un \textbf{método computacional} basado en procesamiento de lenguaje natural (\textit{Sentence-Transformers}) y algoritmos de agrupamiento no supervisado (\textit{clustering} K-Means, Jerárquico y DBSCAN) proyectados dentro de un mismo espacio vectorial latente ($\mathbb{R}^d$). Este método permite caracterizar, mapear y cuantificar numéricamente las diferencias de mallas y los grados de alineación respecto a estándares y vacantes de empleo de LinkedIn (Anexo D).

El producto final operacionaliza el pipeline computacional mediante un tablero de control interactivo (\textit{dashboard}) en React. Este sistema está diseñado bajo una filosofía de soporte a la decisión (\textit{Human-in-the-Loop}), entregando a docentes, coordinadores de carrera y comisiones de rediseño curricular evidencia analítica objetiva —a través de mapas semánticos, matrices de distancia vectorial e índices de brechas— para respaldar la toma de decisiones pedagógicas sin reemplazar el juicio evaluativo humano.

\section{Planteamiento del problema}
\label{cap:problemStatement}

\subsection{Tecnologías y Brechas del Estado del Arte}
\label{sec:literature_gap}

El análisis computacional de la alineación entre programas académicos y demandas del mercado laboral ha evolucionado significativamente en su representación textual. Los enfoques iniciales utilizaron modelos léxicos tradicionales como TF-IDF para ponderar palabras clave en descripciones de cursos; sin embargo, estas representaciones no logran capturar relaciones semánticas complejas ni sinonimias al depender estrictamente de frecuencias léxicas \citep{naseem_comprehensive_2021}. Para superar esto, se adoptaron representaciones distribuidas del lenguaje (\textit{word embeddings} como Word2Vec, GloVe y FastText), que ubican las palabras en espacios vectoriales donde la cercanía denota afinidad conceptual \citep{naseem_comprehensive_2021}. La evidencia empírica muestra que los \textit{word embeddings} superan a la bolsa de palabras, logrando medidas F de alineación mayores (0.571 frente a 0.333) \citep{miura_automated_2021}, y que clasificadores supervisados como SVM alcanzan exactitudes sobresalientes (85.35\%) en la extracción de capacidades \citep{fan_identifying_2023}.

En la actualidad, los modelos de lenguaje basados en transformadores multilingües (\textit{Sentence-Transformers} como SBERT/MPNet) permiten proyectar de manera conjunta textos en español e inglés dentro de un mismo espacio vectorial latente continuo $\mathbb{R}^d$, resolviendo las barreras idiomáticas y capturando la semántica contextual profunda \citep{cristea_course_2025}. Para medir la alineación en este espacio vectorial compartido, la similitud del coseno ($\cos \theta$) y las tasas de cobertura se han consolidado como métricas cuantitativas estándar \citep{miura_automated_2021}. Asimismo, el uso coordinado de algoritmos de agrupamiento no supervisado (\textit{clustering} K-Means, agrupamiento jerárquico y DBSCAN) permite descubrir la estructura disciplinar implícita y aislar asignaturas atípicas sin requerir etiquetado previo, evaluando la calidad geométrica mediante las métricas de Silhouette y Davies-Bouldin \citep{remadi_prompt_2024,fan_identifying_2023}.

Respecto a los antecedentes de la literatura comparada (Anexo B), destaca el trabajo de Russell et al. \citep{russell_e_mapping_2024}, quienes desarrollaron un mapeo computacional de planes de estudio hacia ocupaciones de la industria. No obstante, la solución propuesta en este escrito presenta diferencias sustanciales frente al trabajo de Russell et al.:
\begin{enumerate}[label=\alph*)]
	\item \textbf{Cobertura y Especificidad del Corpus:} Mientras Russell et al. limitaron su análisis a 2 instituciones de EE.UU. en el área acotada de Ciencia de Datos, este trabajo aborda 16 programas universitarios de educación superior en Chile (747 asignaturas), considerando el contexto multilingüe en español y la dualidad de mallas de ciclo largo (Civil) y ejecutoras.
	\item \textbf{Profundidad de Representación Contextual:} Russell et al. emplearon reglas léxicas y ponderación TF-IDF sobre descripciones breves de cursos (SkillNER). En contraste, nuestro método aplica transformadores vectoriales contextuales (SBERT) sobre mallas curriculares y perfiles de egreso declarados.
	\item \textbf{Triangulación Tripartita Simultánea:} A diferencia del mapeo unidireccional de Russell et al. hacia la base O*NET, nuestro método realiza una triangulación simultánea en el mismo espacio vectorial: \textit{Mallas Curriculares} $\leftrightarrow$ \textit{Estándares ACM/IEEE} $\leftrightarrow$ \textit{Ofertas de LinkedIn}.
\end{enumerate}

\subsection{Brecha Tecnológica}
\label{sec:technology_gap}

En el ecosistema universitario chileno no existe una plataforma automatizada implantada que evalúe y compare objetivamente los programas de estudio. Los rediseños curriculares a nivel local dependen de extensas reuniones de comisiones académicas manuales \citep{contreras_diseno_nodate} o de evaluaciones puntuales por pares evaluadores externos en procesos de acreditación \citep{diaz-quezada_rediseno_2019}. Por su parte, la medición de brechas de pertinencia laboral se limita a encuestas de percepción cualitativas aplicadas a egresados y empleadores, propensas al sesgo subjetivo humano \citep{latuz_pina_2024}. A nivel metodológico, las soluciones existentes en la literatura suelen corresponder a flujos experimentales diseñados ad hoc, basados en diccionarios estáticos sin profundidad semántica \citep{constantinov_preliminary_2015} o en modelos que demandan un costoso etiquetado manual de datos \citep{fan_identifying_2023}. Ninguna herramienta integra en una misma arquitectura en español la ingesta de mallas curriculares, el mapeo continuo contra la ACM y el contraste estratificado con el mercado laboral en un dashboard interactivo.

\subsection{Enunciado y Criterios del Problema}
\label{sec:problem_statement}

A partir de la brecha identificada, se formula la siguiente \textbf{pregunta de investigación}:
\begin{quote}
	\textit{“¿En qué medida difieren cuantitativamente la distribución disciplinar y la estructura de las mallas curriculares de Ingeniería Informática en Chile, y cuál es el grado de alineación semántica de sus asignaturas y perfiles de egreso declarados respecto a las directrices de los estándares internacionales ACM/IEEE (CS2013 y CC2020) y a las competencias tecnológicas demandadas por el mercado laboral local extraídas de LinkedIn?”}
\end{quote}

Se concibe como \textbf{conjetura o hipótesis de trabajo} que las mallas curriculares de corte tradicional de ciclo largo (Ingeniería Civil Informática) presentan una alta alineación semántica con los contenidos fundamentales del estándar ACM/IEEE pero manifiestan brechas de cobertura frente a competencias aplicadas de alta demanda en LinkedIn, mientras que las mallas curriculares de corte tecnológico o de ejecución muestran una mayor cercanía vectorial con las vacantes del mercado laboral local pero exhiben menor densidad en las bases teóricas del estándar internacional.

Para abordar esta problemática con rigor técnico y delimitación temporal, el estudio establece los siguientes criterios metodológicos y operacionales:
\begin{enumerate}[label=\arabic*.]
	\item \textbf{Delimitación de las Mallas del Corpus:} Se incluyen programas de formación universitaria vigentes y dictados oficialmente en Chile en el período 2018–2026, considerando tanto mallas de ciclo largo de 5.5 a 6 años como mallas de ejecución de 4 a 4.5 años (Anexo C).
	\item \textbf{Definición de "Diferencias Estructurales de Mallas":} El análisis de variaciones estructurales se evalúa cuantitativamente a través de dos ejes principales: (a) el \textit{Eje de Cobertura Disciplinar} (proporción de créditos y asignaturas de la malla asociados a las distintas \textit{Knowledge Areas} de la ACM), y (b) el \textit{Eje de Balance Modular} (relación entre el núcleo disciplinar obligatorio, líneas de especialización y materias de formación general).
	\item \textbf{Mecanismo de Relación y Estratificación con LinkedIn:} La demanda laboral se modela vectorizando las secciones de requisitos y funciones de 200 vacantes publicadas en LinkedIn durante 2024–2026. Para analizar el mercado de forma representativa, la muestra se divide formalmente en dos estratos: empresas del \textit{Sector Tecnológico (Mundo TI)} y empresas del \textit{Sector No-TI} (industria usuaria: banca, retail, minería, salud). La alineación se calcula evaluando la similitud del coseno $\cos(\mathbf{v}_{\text{asig}}, \mathbf{v}_{\text{job}})$ entre los vectores de las asignaturas de la malla y los requerimientos del mercado.
	\item \textbf{Criterios Numéricos de Decisión:} Se establecen como umbrales de validez: i) \textit{Alineación Normativa Aceptable}: Similitud de coseno promedio $\ge 0.70$ con las \textit{Knowledge Areas} de la ACM y Pureza de agrupamiento $\ge 0.70$; y ii) \textit{Presencia de Brecha Laboral Crítica}: Distancia semántica respecto a las competencias de LinkedIn $> 0.40$ (equivalente a similitud $\cos \theta < 0.60$).
\end{enumerate}
\section{Objetivos del proyecto}
\label{cap:Proposal}

En esta sección se presentan los objetivos que guiarán el desarrollo de este trabajo. En correspondencia directa con la pregunta de investigación planteada en la Sección~\ref{sec:problem_statement}, se formula el objetivo general orientado a definir el alcance técnico y metodológico del proyecto. Posteriormente, se detallan los objetivos específicos que estructuran de forma secuencial y verificable las fases del desarrollo del método computacional, la preparación de los datos, el modelamiento en el espacio vectorial compartida, la validación híbrida y la transferencia de resultados.

\subsection{Objetivo general}

Desarrollar y evaluar un \textbf{método computacional} basado en procesamiento de lenguaje natural, representación semántica vectorial y algoritmos de agrupamiento no supervisado (\textit{clustering}) para analizar, agrupar y comparar cuantitativamente el grado de alineación curricular de las mallas curriculares y perfiles de egreso declarados de las carreras de Ingeniería Informática en Chile con respecto a los estándares de referencia internacional (ACM/IEEE CS2013 y CC2020) y a las competencias demandadas por el mercado laboral tecnológico nacional (estratificado entre el sector TI e industrias no tecnológicas), cuantificando las diferencias estructurales y brechas formativas mediante métricas de similitud semántica en un espacio latente común $\mathbb{R}^d$, y validando los agrupamientos obtenidos mediante un marco de evaluación híbrido (cuantitativo y cualitativo).

\subsection{Objetivos específicos}

\begin{enumerate}[label=\arabic*., nosep]
	\item Construir un corpus curricular y de mercado estructurado, estratificado y normalizado, con el propósito de consolidar una base de datos de texto plano depurada y comparable que sirva de insumo para el análisis semántico de los programas informáticos en Chile, mediante la recopilación curada de mallas curriculares (747 asignaturas) y perfiles de egreso declarados de 16 programas universitarios chilenos, la estructuración digital de las directrices internacionales ACM/IEEE (CS2013 y CC2020), la extracción y estratificación de 200 ofertas laborales de LinkedIn (divididas entre el sector tecnológico especializado y empresas no-TI), el curado manual de etiquetas de referencia (\textit{ground truth}) clasificando las 747 asignaturas en las 18 \textit{Knowledge Areas} de la ACM, y la aplicación de técnicas avanzadas de preprocesamiento NLP (remoción de \textit{stop words} académicas, normalización ortográfica y lematización). \textbf{Criterio de calidad}: Cobertura comprobada del 100\% de las mallas y asignaturas del corpus muestral (747 asignaturas en 16 programas) con una tasa de datos corruptos o nulos $< 1.0\%$ tras la etapa de normalización y filtrado de ruido.

	\item Generar representaciones vectoriales densas (\textit{embeddings}) de los objetos curriculares y laborales en un espacio latente común ($\mathbb{R}^d$), a fin de mapear las asignaturas de las mallas, los perfiles declarados, los estándares ACM y las vacantes de mercado en una misma geometría continua que represente su semántica contextual y habilite el cálculo matemático de similitudes, utilizando modelos avanzados basados en transformadores multilingües (\textit{Sentence-Transformers} SBERT/MPNet) en combinación con algoritmos de reducción de dimensionalidad (como UMAP o t-SNE) para su visualización espacial. \textbf{Criterio de calidad}: Conservación de la estructura topológica local del espacio vectorial verificada mediante una métrica de fidelidad de proyección (\textit{Trustworthiness} $\ge 0.85$).

	\item Identificar perfiles formativos específicos y medir la alineación semántica y las diferencias de mallas, con el fin de detectar patrones de especialización implícitos entre las distintas universidades y cuantificar con rigor matemático la correspondencia y las brechas curriculares frente a las directrices de referencia y las necesidades de la industria, mediante la aplicación coordinada de una estrategia tripartita de agrupamiento no supervisado (K-Means para particiones convexas con centroides explícitos, Agrupamiento Jerárquico para dendrogramas de parentesco, y DBSCAN para el aislamiento de asignaturas atípicas u \textit{outliers}), identificando patrones estructurales \textit{(e.g., concentración temprana de ciencias básicas vs. especialización precoz; profundidad de la línea de programación; presencia de ciberseguridad o inteligencia artificial en la malla)} y perfiles formativos \textit{(e.g., Malla orientada a Ingeniería de Software, Malla enfocada en Ciencia de Datos, Malla de Infraestructura y Redes, Malla de Computación Teórica)}, junto al cálculo de similitud del coseno ($\cos \theta$), tasa de cobertura ($\% \text{Cobertura}$) e índice de brecha semántica ($1 - \cos \theta$). \textbf{Criterio de calidad}: Obtención de una partición coherente con un Coeficiente de Silhouette $\ge 0.25$ y un Índice de Davies-Bouldin $< 1.50$.

	\item Validar la calidad, estabilidad y aplicabilidad práctica de los agrupamientos y brechas obtenidas, a objeto de asegurar la consistencia técnica del método computacional y la robustez de los hallazgos para su uso en la gestión académica, a través de un marco de evaluación híbrido que combine métricas cuantitativas intrínsecas, validación externa contrastada contra la taxonomía ACM utilizando el \textit{ground truth} de 747 asignaturas etiquetadas mediante el Índice Rand Ajustado ($\text{ARI} \ge 0.35$), Información Mutua Normalizada ($\text{NMI} \ge 0.40$) y Pureza ($\text{Purity} \ge 0.70$), evaluación de estabilidad mediante re-muestreo bootstrap ($\text{ARI}_{\text{promedio}} \ge 0.75$) y juicio cualitativo por un panel de 3 expertos independientes evaluado con un coeficiente de acuerdo inter-evaluador de Fleiss/Cohen $\kappa \ge 0.60$. \textbf{Criterio de calidad}: Cumplimiento simultáneo de los umbrales de validación intrínseca, externa, de estabilidad y de concordancia inter-experto estipulados.

	\item Operacionalizar el método computacional y sus resultados analíticos en una interfaz interactiva de apoyo a la toma de decisiones (\textit{Human-in-the-Loop}), para transferir el conocimiento generado a los gestores de las instituciones de educación superior y apoyar las deliberaciones informadas en los procesos de rediseño curricular continuo, mediante la implementación de un tablero de control (\textit{dashboard}) interactivo en React que integre mapas semánticos de asignaturas de las mallas, matrices de calor de competencias y gráficos dinámicos de brechas formativas. \textbf{Criterio de calidad}: Despliegue funcional de un tablero de control navegable y responsivo con una latencia media de respuesta $< 2.0$ segundos por consulta interactiva y usabilidad validada sin errores críticos de renderizado.
\end{enumerate}

% Entrega 2
\section{Análisis del enfoque}
\label{cap:approachAnalysis}

\subsection{Representación Vectorial Contextual en Espacio Compartido}
\label{sec:nlp_representation}

Para vectorizar los textos del corpus sin sufrir las limitaciones de los modelos de bolsa de palabras (TF-IDF) ni la falta de contexto de los \textit{word embeddings} estáticos, se seleccionan modelos de lenguaje basados en transformadores multilingües (\textit{Sentence-Transformers} como \textit{paraphrase-multilingual-mpnet-base-v2} o \textit{multilingual-e5-large}). Estos modelos permiten proyectar los títulos y descriptores modulares de las asignaturas de las mallas, las oraciones de los perfiles de egreso declarados, las definiciones de las \textit{Knowledge Areas} de la ACM/IEEE y las secciones de requerimientos de las vacantes de LinkedIn dentro de un \textbf{mismo espacio vectorial latente continuo ($\mathbb{R}^d$)}. Esta proyección homotópica habilita el cálculo directo de distancia euclidiana y similitud de coseno ($\cos \theta$) entre elementos de distintas fuentes de datos en español e inglés.

La alineación y las brechas semánticas entre una malla curricular $M_i$ y una fuente de referencia (ACM o LinkedIn) se cuantifican mediante tres métricas principales:
\begin{enumerate}[label=\alph*)]
	\item \textbf{Similitud de Coseno ($\cos \theta$):} Mide la orientación relativa entre el vector de una asignatura $\mathbf{v}_{\text{asig}}$ y el vector de una competencia o área de conocimiento $\mathbf{v}_{\text{ref}}$:
	      \begin{equation}
		      \cos(\mathbf{v}_{\text{asig}}, \mathbf{v}_{\text{ref}}) = \frac{\mathbf{v}_{\text{asig}} \cdot \mathbf{v}_{\text{ref}}}{\|\mathbf{v}_{\text{asig}}\| \|\mathbf{v}_{\text{ref}}\|}
	      \end{equation}
	\item \textbf{Tasa de Cobertura de la Malla ($\% \text{Cobertura}$):} Representa el porcentaje de conceptos clave o \textit{Knowledge Areas} del estándar internacional o del mercado laboral representados en la malla curricular con una similitud semántica $\cos \theta \ge 0.60$.
	\item \textbf{Índice de Brecha Semántica ($1 - \cos \theta$):} Cuantifica la distancia vectorial que refleja el déficit o la ausencia explícita de asignaturas en la malla curricular para cubrir requerimientos tecnológicos específicos.
\end{enumerate}

\subsection{Justificación de la Estrategia Tripartita de Agrupamiento}
\label{sec:clustering_justification}

Para descubrir la estructura disciplinar de las mallas curriculares y agrupar las asignaturas de forma no supervisada, se adopta una \textbf{estrategia comparativa tripartita} que combina K-Means, Agrupamiento Jerárquico (\textit{Hierarchical Clustering}) y DBSCAN. Esta combinación no es arbitraria; responde a que cada algoritmo resuelve una necesidad analítica complementaria e insustituible sobre los datos curriculares:
\begin{enumerate}[label=\arabic*.]
	\item \textbf{K-Means Clustering:} Produce particiones convexas en $k$ grupos maximizando la homogeneidad interna. Su principal ventaja matemática radica en que calcula centroides explícitos $\mathbf{c}_k \in \mathbb{R}^d$ para cada clúster, los cuales actúan como la representación vectorial promedio de un área disciplinar dentro de la malla. El número óptimo de clústeres $k$ se determina mediante la maximización del Coeficiente de Silhouette.
	\item \textbf{Agrupamiento Jerárquico Aglomerativo (Ward Linkage):} Construye una estructura arborescente o dendrograma que preserva las relaciones de jerarquía y subordinación taxonómica entre asignaturas (e.g., relacionando asignaturas de programación básica con materias de ingeniería de software avanzada). Permite visualizar los niveles de abstracción disciplinar sin prefijar el número de grupos.
	\item \textbf{DBSCAN (Density-Based Spatial Clustering of Applications with Noise):} Algoritmo basado en densidad que define clústeres como regiones densas de vectores en $\mathbb{R}^d$ separadas por zonas de baja densidad. Su función técnica crucial en este proyecto es la \textbf{detección y aislamiento explícito de asignaturas atípicas o ruido (\textit{outliers})}, es decir, asignaturas singulares presentes en mallas de ciertas universidades que no encajan en las agrupaciones estándar del sistema.
\end{enumerate}

Para evaluar la validez externa de los clústeres generados frente a un estándar objetivo, se realiza una anotación manual de las 747 asignaturas del corpus clasificándolas en las 18 \textit{Knowledge Areas} de la ACM (CS2013). Esta clasificación actúa como etiqueta de referencia (\textit{ground truth}) para calcular el Índice Rand Ajustado ($\text{ARI} \ge 0.35$), la Información Mutua Normalizada ($\text{NMI} \ge 0.40$) y la Pureza del agrupamiento ($\text{Purity} \ge 0.70$).

\section{Descripción del producto}
\label{cap:productDescription}

Esta sección detalla la arquitectura del producto desarrollado, el cual comprende el pipeline del método computacional de análisis semántico y el tablero de control interactivo (\textit{dashboard}) para la comparación visual de mallas curriculares.

\subsection{Flujo Metodológico e Insumos}
\label{sec:flow_inputs}

El producto se estructura en un pipeline computacional compuesto por cinco fases correlativas (detalladas en el Anexo E):
\begin{enumerate}[label=\arabic*.]
	\item \textbf{Fase 1: Adquisición y Preprocesamiento:} Ingesta de mallas curriculares, perfiles de egreso declarados, estándares y vacantes; limpieza de texto, remoción de \textit{stop words} académicas y lematización.
	\item \textbf{Fase 2: Generación de Embeddings Contextuales:} Vectorización multilingüe mediante \textit{Sentence-Transformers} (SBERT/MPNet) proyectando los textos en un espacio vectorial común $\mathbb{R}^d$, y reducción dimensional con UMAP/t-SNE.
	\item \textbf{Fase 3: Modelamiento y Agrupamiento Semántico:} Ejecución de la estrategia tripartita de \textit{clustering} (K-Means, Jerárquico y DBSCAN) para la identificación de patrones estructurales y asignaturas atípicas.
	\item \textbf{Fase 4: Comparación y Cálculo de Brechas:} Generación de matrices de similitud de coseno ($\cos \theta$), tasas de cobertura ($\% \text{Cobertura}$) e índices de brecha semántica.
	\item \textbf{Fase 5: Despliegue en Dashboard React:} Exportación de matrices en formato JSON y renderización de mapas semánticos e indicadores interactivos.
\end{enumerate}

La recolección de insumos contempla tres dimensiones rigurosamente delimitadas:
\begin{itemize}
	\item \textbf{Dimensión Académica Nacional:} 16 mallas curriculares (747 asignaturas) e itinerarios formativos vigentes (2018–2026), junto a los perfiles de egreso declarados por las universidades (Anexo C).
	\item \textbf{Dimensión Normativa Internacional:} Las 18 \textit{Knowledge Areas} y taxonomías de los estándares ACM/IEEE CS2013 y CC2020 \citep{cc2020_task_force_computing_2020}.
	\item \textbf{Dimensión Mercado Laboral:} $N = 200$ convocatorias de empleo recolectadas de LinkedIn (2024–2026), estratificadas formalmente entre empresas del \textit{Sector Tecnológico (Mundo TI)} y empresas del \textit{Sector No-TI} (banca, minería, retail, servicios públicos) (Anexo D).
\end{itemize}

\subsection{Algoritmo del Método Computacional}
\label{sec:algorithm_pseudocode}

A continuación se presenta el pseudocódigo formal que describe el funcionamiento secuencial del método computacional propuesto:

\vspace{0.3cm}

\begin{algorithm}[H]
	\caption{Pipeline Computacional de Análisis Semántico y Comparación de Mallas Curriculares}
	\label{alg:pipeline_mallas_doc}
	\small
	\begin{algorithmic}[1]
		\Require Corpus de mallas $\mathcal{M} = \{M_1, M_2, \dots, M_n\}$, Perfiles de Egreso $\mathcal{P}$, Estándares $\mathcal{S}$ (ACM/IEEE KAs), Vacantes LinkedIn $\mathcal{L}$, Modelo SBERT $\mathcal{E}$
		\Ensure Representaciones vectoriales, Clústeres de mallas, Matrices de Similitud $S$ y Dashboard en React
		\Statex
		\Procedure{PipelineAnalisisCurricular}{$\mathcal{M}, \mathcal{P}, \mathcal{S}, \mathcal{L}, \mathcal{E}$}
			\State \textbf{Fase 1: Preprocesamiento y Limpieza}
			\For{cada asignatura $a \in M_i \in \mathcal{M}$}
				\State $a_{\text{clean}} \leftarrow \text{NormalizarOrtografia}(\text{RemoverStopWords}(a.\text{nombre} \cup a.\text{descriptor}))$
			\EndFor
			\For{cada $p \in \mathcal{P}$, $s \in \mathcal{S}$, $l \in \mathcal{L}$}
				\State $p_{\text{clean}} \leftarrow \text{LimpiezaTexto}(p)$, \quad $s_{\text{clean}} \leftarrow \text{LimpiezaTexto}(s)$, \quad $l_{\text{clean}} \leftarrow \text{LimpiezaTexto}(l)$
			\EndFor
			
			\Statex
			\State \textbf{Fase 2: Generación de Embeddings Multilingües}
			\State $\mathbf{V}_{\mathcal{M}} \leftarrow \mathcal{E}.\text{encode}(\mathcal{M}_{\text{clean}})$, \quad $\mathbf{V}_{\mathcal{P}} \leftarrow \mathcal{E}.\text{encode}(\mathcal{P}_{\text{clean}})$
			\State $\mathbf{V}_{\mathcal{S}} \leftarrow \mathcal{E}.\text{encode}(\mathcal{S}_{\text{clean}})$, \quad $\mathbf{V}_{\mathcal{L}} \leftarrow \mathcal{E}.\text{encode}(\mathcal{L}_{\text{clean}})$
			
			\Statex
			\State \textbf{Fase 3: Modelamiento y Agrupamiento Semántico (Clustering)}
			\State $\mathcal{C}_{\text{kmeans}} \leftarrow \text{KMeans}(\mathbf{V}_{\mathcal{M}}, k)$, \quad $\mathcal{C}_{\text{hier}} \leftarrow \text{HierarchicalClustering}(\mathbf{V}_{\mathcal{M}})$
			\State $\mathcal{O}_{\text{dbscan}} \leftarrow \text{DBSCAN}(\mathbf{V}_{\mathcal{M}}, \epsilon, \text{min\_pts})$ \Comment{Aislamiento de asignaturas atípicas}
			
			\Statex
			\State \textbf{Fase 4: Comparación y Alineación Vectorial}
			\For{cada programa $M_i \in \mathcal{M}$}
				\State $S_{\text{ACM}}[i] \leftarrow \text{SimilitudCoseno}(\mathbf{V}_{M_i}, \mathbf{V}_{\mathcal{S}})$ \Comment{Alineación con Estándar ACM/IEEE}
				\State $S_{\text{LinkedIn}}[i] \leftarrow \text{SimilitudCoseno}(\mathbf{V}_{M_i}, \mathbf{V}_{\mathcal{L}})$ \Comment{Alineación con Mercado TI / No-TI}
				\State $S_{\text{Perfil}}[i] \leftarrow \text{SimilitudCoseno}(\mathbf{V}_{M_i}, \mathbf{V}_{P_i})$ \Comment{Alineación Malla vs Perfil}
			\EndFor
			
			\Statex
			\State \textbf{Fase 5: Despliegue en Dashboard React}
			\State Exportar matrices $\{S_{\text{ACM}}, S_{\text{LinkedIn}}, S_{\text{Perfil}}\}$ y coordenadas UMAP a formato JSON
			\State Renderizar mapas semánticos interactivos y matrices de diferencias de mallas en React
		\EndProcedure
	\end{algorithmic}
\end{algorithm}

\vspace{0.3cm}

\subsection{Marco de Evaluación Híbrido}
\label{sec:evaluation_framework}

El marco de validación se compone de cuatro componentes complementarios:
\begin{enumerate}[label=\alph*)]
	\item \textbf{Evaluación Cuantitativa Intrínseca y Externa:} Validación geométrica con Silhouette ($\ge 0.25$) y Davies-Bouldin ($< 1.50$), y validación externa contra las 18 áreas del estándar ACM mediante el curado manual de las 747 asignaturas como etiqueta de referencia (\textit{ground truth}), exigiendo $\text{ARI} \ge 0.35$, $\text{NMI} \ge 0.40$ y $\text{Purity} \ge 0.70$.
	\item \textbf{Análisis de Estabilidad por Bootstrap:} Evaluación de consistencia algorítmica ($B = 100$ iteraciones sobre subconjuntos aleatorios del 80\%–90\% del corpus) asegurando $\text{ARI}_{\text{estabilidad}} \ge 0.75$.
	\item \textbf{Evaluación Cualitativa por Panel de Expertos:} Revisión por un panel independiente de 3 expertos académicos/industriales evaluando la coherencia pedagógica con una concordancia inter-evaluador $\kappa \ge 0.60$.
	\item \textbf{Evaluación de Usabilidad y Despliegue del Dashboard:} Verificación de latencia interactiva $< 2.0$ segundos en la plataforma React para garantizar su efectividad en reuniones de comisiones académicas.
\end{enumerate}


% Entrega 3
\section{Metodología, herramientas y ambiente de desarrollo}
\label{cap:Methodology}

\subsection{Justificación de la metodología a utilizar}
\label{sec:crisp_dm}
Para estructurar y guiar el desarrollo de este proyecto, se ha seleccionado la metodología CRISP-DM (\textit{Cross-Industry Standard Process for Data Mining}). La elección de este marco de trabajo se justifica por su enfoque iterativo y su probada eficacia en proyectos de ciencia de datos y minería de texto. CRISP-DM proporciona un ciclo de vida claro que abarca desde la comprensión del dominio educativo e industrial hasta el despliegue del \textit{dashboard} interactivo, asegurando que cada fase del procesamiento de lenguaje natural y modelamiento no supervisado responda de manera directa a los objetivos y necesidades del problema planteado.

\subsection{Herramientas de desarrollo}
\label{sec:tools}
El procesamiento analítico y la implementación tecnológica se llevarán a cabo utilizando el lenguaje de programación Python, debido a su robusto ecosistema orientado al procesamiento de lenguaje natural y aprendizaje automático. Entre las principales librerías a utilizar destacan:
\begin{itemize}
    \item \textbf{Procesamiento de Lenguaje Natural (NLP)}: \textit{spaCy} y \textit{NLTK} para las tareas de limpieza, tokenización, remoción de palabras vacías y lematización.
    \item \textbf{Representación Vectorial}: \textit{Sentence-Transformers} (basado en la biblioteca \textit{Hugging Face Transformers}) para la generación de \textit{embeddings} semánticos.
    \item \textbf{Agrupamiento y Reducción de Dimensionalidad}: \textit{scikit-learn} para algoritmos como K-Means, DBSCAN y cálculo de similitud del coseno; junto con librerías especializadas como \textit{UMAP} y \textit{t-SNE} para la proyección espacial de los datos.
    \item \textbf{Visualización y Despliegue}: \textit{Matplotlib} y \textit{Seaborn} para el análisis exploratorio, y \textit{React} para la construcción de la interfaz gráfica y el tablero interactivo final.
\end{itemize}

\subsection{Ambiente de desarrollo}
\label{sec:environment}
El entorno de trabajo estará centralizado de forma nativa en el sistema operativo \textit{Ubuntu 24.04 LTS}. Todo el desarrollo del código fuente, el diseño del \textit{pipeline} y el empaquetado de la solución se realizarán utilizando \textit{Visual Studio Code} (VS Code) como entorno de desarrollo integrado (IDE). Además, para garantizar la trazabilidad íntegra del proyecto, la totalidad de los archivos (tanto el código como la documentación) quedarán estrictamente versionados utilizando el sistema \textit{Git}, con el repositorio central alojado en la plataforma \textit{GitHub}.

\section{Plan de trabajo}
\label{cap:Schedule}

\subsection{Definición de Actividades e Hitos}
\label{sec:activities_milestones}

Las actividades del proyecto se estructuran temporalmente a lo largo de un semestre académico (16 semanas de trabajo), organizadas según las fases de la metodología CRISP-DM y vinculadas de forma directa con los Objetivos Específicos (OE1 a OE5).

Para garantizar la verificabilidad del avance del proyecto y la entrega oportuna de resultados, se definen cinco \textbf{Hitos de Proyecto (H1 a H5)} con sus correspondientes entregables explícitos y semanas de término, tal como se sintetiza en la Tabla~\ref{tab:hitos_planification}:

\vspace{0.3cm}

\begin{table}[H]
	\centering
	\small
	\renewcommand{\baselinestretch}{1.0}\selectfont
	\caption{Hitos, entregables explícitos y vinculación con Objetivos Específicos.}
	\label{tab:hitos_planification}
	\captionsetup{skip=6pt}
	\begin{tabularx}{\textwidth}{c p{3.2cm} X c c}
		\toprule
		\textbf{Hito} & \textbf{Fase CRISP-DM} & \textbf{Entregable / Resultado Explícito} & \textbf{OE} & \textbf{Semana} \\
		\midrule
		\textbf{H1} & 1. Comprensión & Corpus crudo compilado (16 mallas curriculares, 747 asignaturas, perfiles de egreso declarados, ACM, 200 vacantes de LinkedIn). & OE1 & Sem 2 \\
		\textbf{H2} & 2. Preparación & Dataset de mallas y perfiles estructurado y normalizado en JSON/SQL, y \textit{ground truth} etiquetado de 747 asignaturas. & OE1 & Sem 4 \\
		\textbf{H3} & 3. Modelamiento & Matriz de \textit{embeddings} SBERT de asignaturas de mallas y modelos de \textit{clustering} (K-Means, Jerárquico, DBSCAN). & OE2, OE3 & Sem 8 \\
		\textbf{H4} & 4. Evaluación & Reporte cuantitativo de métricas (Silhouette, Davies-Bouldin, ARI, NMI, Pureza), análisis bootstrap y matriz de diferencias de mallas. & OE3, OE4 & Sem 12 \\
		\textbf{H5} & 5. Despliegue & Dashboard interactivo de comparación de mallas desplegado en React y reporte de validación cualitativa por expertos. & OE4, OE5 & Sem 16 \\
		\bottomrule
	\end{tabularx}
\end{table}

\vspace{0.3cm}

\subsection{Carta Gantt}
\label{sec:gantt_chart}

La planificación temporal detallada de las tareas, dependencias entre fases y asignación de semanas a lo largo del periodo académico se presenta en la representación gráfica de la Carta Gantt ubicada en la Tabla del Anexo F.




\newpage
\addcontentsline{toc}{section}{Bibliografía}
\bibliographystyle{apalike}
\bibliography{References/references.bib}

\newpage
\section*{Anexos}
\addcontentsline{toc}{section}{Anexos}
\setcounter{subsection}{0}
\renewcommand{\thesubsection}{Anexo \Alph{subsection}}
\renewcommand{\theHsubsection}{Anexo.\arabic{subsection}}
\addtocontents{toc}{\string\renewcommand{\string\numberline}{\string\appendixnumberline}}

\subsection{Glosario}

\begin{itemize}
    \item \textbf{Alineación Curricular}: Grado de correspondencia, congruencia y consistencia entre los componentes de un programa educativo (tales como perfiles de egreso, planes de estudio y descriptores de asignaturas) con respecto a directrices internacionales y los requerimientos del entorno laboral.
    
    \item \textbf{Clustering (Agrupamiento)}: Técnica de aprendizaje no supervisado que consiste en agrupar un conjunto de objetos de tal manera que los miembros del mismo grupo (\textit{clúster}) sean más similares entre sí que con los de otros grupos, basándose en alguna métrica de distancia o similitud geométrica.
    
    \item \textbf{Coeficiente de Silueta (\textit{Silhouette Coefficient})}: Métrica de validación interna de algoritmos de \textit{clustering} que evalúa qué tan bien agrupado está un elemento. Mide la cercanía de un objeto a su propio clúster (cohesión) frente a su cercanía a otros clústeres (separación), variando entre -1 y 1.
    
    \item \textbf{Corpus}: Conjunto estructurado y homogéneo de textos escritos, recopilado de manera sistemática para ser utilizado como base de análisis lingüístico o procesamiento computacional.
    
    \item \textbf{DBSCAN (\textit{Density-Based Spatial Clustering of Applications with Noise})}: Algoritmo de agrupamiento no supervisado basado en densidad que agrupa puntos cercanos y etiqueta como valores atípicos (\textit{outliers}) aquellos que se encuentran aislados en regiones de baja densidad.
    
    \item \textbf{Embeddings (Incrustaciones Vectoriales)}: Representaciones matemáticas de textos (palabras, frases o documentos completos) en forma de vectores densos en un espacio continuo de baja dimensionalidad, donde la cercanía geométrica de los vectores representa similitud semántica.
    
    \item \textbf{Índice de Davies-Bouldin}: Métrica de evaluación intrínseca de algoritmos de agrupamiento que calcula la similitud promedio entre cada clúster y su clúster más similar. Valores más bajos indican agrupamientos con clústeres más compactos y mejor separados.
    
    \item \textbf{K-Means}: Algoritmo clásico de agrupamiento no supervisado que particiona un conjunto de datos en $K$ clústeres distintos y disjuntos, asignando cada elemento al clúster cuyo centroide se encuentra a menor distancia geométrica (generalmente distancia euclidiana).
    
    \item \textbf{Lematización}: Proceso de preprocesamiento de lenguaje natural que consiste en reducir las palabras de un texto a su forma básica o raíz morfológica canonizada (lema), como infinitivos para verbos y masculino singular para adjetivos y sustantivos.
    
    \item \textbf{Procesamiento de Lenguaje Natural (PLN / NLP)}: Subcampo de la inteligencia artificial y la lingüística computacional enfocado en el diseño e implementación de algoritmos que permiten a las computadoras procesar, entender y analizar textos en lenguaje humano.
    
    \item \textbf{Sentence-Transformers (SBERT)}: Framework de aprendizaje profundo basado en redes neuronales siamesas, utilizado para adaptar modelos de lenguaje tipo BERT para generar embeddings de oraciones y textos completos de alta calidad semántica.
    
    \item \textbf{Similitud del Coseno}: Métrica matemática utilizada para cuantificar la similitud entre dos vectores en un espacio multidimensional, midiendo el coseno del ángulo formado entre ellos, independientemente de su magnitud.
    
    \item \textbf{Stop Words (Palabras Vacías)}: Términos que carecen de significado informativo relevante por sí mismos (como artículos, preposiciones y conjunciones, o en el contexto académico, términos como ``asignatura'', ``horas'' o ``carrera''), que son filtrados durante el preprocesamiento de texto.
    
    \item \textbf{TF-IDF (\textit{Term Frequency - Inverse Document Frequency})}: Técnica estadística clásica de representación de textos que evalúa la importancia de una palabra en un documento en relación con un corpus completo, ponderándola según su frecuencia local y su rareza global.
    
    \item \textbf{t-SNE (\textit{t-Distributed Stochastic Neighbor Embedding})}: Algoritmo no lineal de reducción de dimensionalidad utilizado principalmente para la visualización de datos de alta dimensionalidad en espacios de dos o tres dimensiones.
    
    \item \textbf{UMAP (\textit{Uniform Manifold Approximation and Projection})}: Algoritmo moderno de reducción de dimensionalidad no lineal fundamentado en la geometría de variedades, eficiente computacionalmente y robusto para preservar la estructura local y global de los datos.
\end{itemize}

%Results from the state-of-the-art strategies
%\setcounter{page}{1}
\subsection{Resumen comparativo de trabajos relacionados sobre alineación curricular automatizada.}
\begin{landscape}
\footnotesize
\renewcommand{\baselinestretch}{1.0}\selectfont
\renewcommand{\arraystretch}{1.2}
\begin{xltabular}{\linewidth}{p{2.2cm} X X X X}
    \caption{Análisis comparativo de objetivos, metodologías, resultados y limitaciones de la literatura de mapeo curricular.} \label{tab:estado_arte} \\
    \toprule
    \textbf{Referencia} &
    \textbf{Objetivo}   &
    \textbf{Metodología} &
    \textbf{Resultados principales} &
    \textbf{Limitaciones identificadas} \\
    \midrule
    \endhead

    \citep{russell_e_mapping_2024} & 
    Automatizar el mapeo de planes de estudio hacia habilidades y ocupaciones de la industria utilizando descripciones de cursos. & 
    Procesamiento de Lenguaje Natural (NLP) usando SkillNER, extracción de características con ponderación TF-IDF y similitud del coseno para emparejar con O*NET. & 
    Extracción promedio de 5.3 a 10.4 habilidades por curso. Fuerte precisión (precision@5 del 60\% al 100\%). Fuerte señal de habilidades incluso en textos breves. & 
    Evaluación limitada a un campo (ciencia de datos) y dos instituciones. Falsos positivos por coincidencia léxica. Falta información sobre la profundidad de la habilidad enseñada. \\
    \addlinespace
    
    \citep{cristea_course_2025} & 
    Evaluar el rendimiento de LLM frente a métodos NLP tradicionales en la extracción de habilidades para la analítica curricular. & 
    Comparación de estrategias: Generación Aumentada por Recuperación (RAG), prompting Zero-shot (GPT, Llama, Gemini, Claude), modelos basados en embeddings (SBERT) y TF-IDF en 400 documentos. & 
    RAG fue la estrategia con mejor desempeño en todos los documentos. El prompting Zero-shot tuvo peor rendimiento que métodos NLP tradicionales en la mayoría de los casos. & 
    El rendimiento varía significativamente según el modelo seleccionado y las estrategias de prompting. Requiere diseño y ajuste cuidado antes de implementaciones a gran escala. \\
    \addlinespace
    
    \citep{cristea_artificial_2025} & 
    Evaluar el rol de la IA Generativa (LLMs) como co-revisor curricular para generar mapeos ponderados de habilidades en evaluaciones. & 
    Uso de GPT-4.0 frente a métodos ML (BERT, TF-IDF) para mapear 7 habilidades empresariales. Uso del alfa de Krippendorff para comparar resultados con el juicio de expertos. & 
    Los mapeos generados por LLM se alinearon más con el juicio experto (alfa 0.76) que los métodos ML (0.65). Alta precisión en competencias como Comunicación Oral y Trabajo en Equipo. & 
    Limitado a un solo LLM y programa de estudios. Desempeño débil en habilidades que se infieren del contexto (ej. Comunicación escrita) y no se detallan en el texto. \\
    \addlinespace
    
    \citep{fan_identifying_2023} & 
    Identificar capacidades digitales evaluadas en rúbricas de cursos universitarios de manera automatizada utilizando Machine Learning. & 
    Extracción de características mediante TF-IDF. Comparación de clasificadores: Redes Neuronales (ANN), Máquinas de Vectores de Soporte (SVM), Random Forest y XGBoost. & 
    SVM obtuvo el mejor rendimiento (exactitud 0.8535 y F1-score 0.8338). El rendimiento mejora y luego decae al aumentar las características extraídas. & 
    Dependencia de la intervención humana para etiquetar los datos de entrenamiento (solo 1.735 muestras etiquetadas). Los algoritmos probados no pueden utilizar pesos difusos para mitigar ruido. \\
    \addlinespace
    
    \citep{duarte_broader_2023} & 
    Generar representaciones visuales comprensibles de los planes de estudio para reducir brechas de comunicación entre los actores. & 
    Clasificación usando índices de Wikipedia, NLP (tokenización y normalización), y visualización mediante ideogramas circulares y grafos de red. & 
    El mapeo apoyó con evidencia visual las debilidades del programa. Las representaciones cuantitativas empoderan el entendimiento entre diferentes grupos de interés. & 
    Las representaciones agrupan conceptos pero no trazan rutas de aprendizaje precisas o de precedencia detalladas, requiriendo revisión y otras herramientas para detalles exactos. \\
    \addlinespace
    
    \citep{miura_automated_2021} & 
    Mapear automáticamente programas educativos de ciencias de la computación frente a estándares de aprendizaje (ITSS+) para reducir costos manuales. & 
    Representación de texto usando Bag-of-Words (BoW) y Word Embeddings (Word2Vec). Emparejamiento usando similitud del coseno y algoritmos de aprendizaje supervisado. & 
    El aprendizaje supervisado con Word Embeddings alcanzó la mayor medida F (0.571), superando ampliamente a enfoques tradicionales como BoW (0.333) al capturar semántica abstracta. & 
    Requiere aumentar los datos de entrenamiento y cuantificar más correctamente características ponderando partes importantes del material didáctico. \\
    \addlinespace
    
    \citep{constantinov_preliminary_2015} & 
    Ajustar un currículo de ingeniería extrayendo las demandas del mercado laboral a través de redes profesionales. & 
    Pipeline de datos para Análisis de Redes Profesionales (PNA) con perfiles de LinkedIn. Extracción de habilidades en sílabos mediante TF-IDF y diccionarios personalizados. & 
    El sistema identificó habilidades del plan de estudios obsoletas en la industria y habilidades muy demandadas en el trabajo que el programa académico no estaba cubriendo. & 
    Metodología basada en diccionarios personalizados estáticos y extracción mediante palabras clave sin abordajes semánticos complejos o contextuales de modelado de lenguaje. \\
  
    \bottomrule
\end{xltabular}
\end{landscape}




% Anexo 3
\subsection{Descripción del corpus de datos}
\label{app:corpus}
 \begin{table}[H]
        \centering
        \small
        % interlineado simple dentro de la tabla (1.0)
        \renewcommand{\baselinestretch}{1.0}\selectfont
        \caption{Planes de estudio y asignaturas recopiladas por universidad.}
        \label{tab:corpus_universidades}
        \captionsetup{skip=8pt}                                                                                              
                                                                                                                             
        % Definimos el contador de filas (se resetea al inicio de la tabla)
        \newcounter{contadorFilas}
        \setcounter{contadorFilas}{0}
        
        % Comando para terminar la fila incrementando el contador y dibujando una línea horizontal
        \newcommand{\finFila}{%
          \stepcounter{contadorFilas}%
          \\ \hline%
        }
        
        % Definición de columnas con alineación a la izquierda (raggedright) para evitar espacios gigantes
        \begin{tabularx}{\textwidth}{>{\raggedright\arraybackslash}X >{\raggedright\arraybackslash}X c}                                                                              
            \toprule                                                                                                         
            \textbf{Universidad} &                                                                                           
            \textbf{Carrera / Programa} &                                                                                    
            \textbf{N\textsuperscript{o} Asignaturas} \\                                                                                    
            \midrule                                                                                                         
                                                                                                                             
            % Fila 1                                                                                              
            Universidad de Santiago de Chile &                                                                               
            Ingeniería de Ejecución en Computación e Informática &                                                                                
            43 \finFila                                                                                       
                                                                                                                             
            % Fila 2                                                                                              
            Universidad de Santiago de Chile &                                                                                           
            Ingeniería Civil en Informática &                                                                                
            58 \finFila
            
            % Fila 3 
            Universidad de Chile & 
            Ingeniería Civil en Computación & 
            54 \finFila
            
            % Fila 4
            Universidad Andrés bello & 
            Ingeniería en Computación e informática & 
            41 \finFila
            
            % Fila 5
            Universidad de Magallanes & 
            Ingeniería en Computación e Informática & 
            41 \finFila

            % Fila 6
            Universidad Tecnológica Metropolitana & 
            Ingeniería en Informática & 
            44 \finFila

            % Fila 7
            Universidad de Concepción & 
            Ingeniería Civil Informática & 
            54 \finFila

            % Fila 8
            Universidad Técnica Federico Santa María & 
            Ingeniería en Informática & 
            42 \finFila

            % Fila 9
            Pontificia Universidad Catolica de Valparaiso  & 
            Ingeniería Informática & 
            44 \finFila

            % Fila 10
            DuocUC  & 
            Ingeniería Informática & 
            33 \finFila

            % Fila 11
            Universidad San Sebastián  & 
            Ingeniería Civil en Informática & 
            50 \finFila

            % Fila 12
            Universidad Santo Tomás   & 
            Ingeniería Civil en Informatica y Sistemas Inteligentes  & 
            60 \finFila

            % Fila 13
            Universidad de la Frontera   & 
            Ingeniería Civil Informática  & 
            57 \finFila

            % Fila 14
            Universidad del Bío-Bío   & 
            Ingeniería de Ejecución en Computación e Informática  & 
            46 \finFila

            % Fila 15
            Universidad de Talca   & 
            Ingeniería Civil en Computación  & 
            43 \finFila

            % Fila 16
            Universidad Autonoma  & 
            Ingeniería en Informática  & 
            32 \finFila
            
            \multicolumn{2}{l}{\textbf{Total del Corpus (16 Programas / Planes de Estudio)}} & \textbf{747} \\
            \bottomrule
        \end{tabularx}
    \end{table}

% Anexo 4
\subsection{Protocolo y caracterización de la muestra de datos de LinkedIn}
\label{app:linkedin_sample}

Con el fin de garantizar la representatividad, validez y reproducibilidad del análisis de pertinencia laboral, se define un protocolo formal para la selección y curaduría de la muestra de ofertas de empleo provenientes de la plataforma LinkedIn:

\begin{enumerate}[label=\alph*.]
    \item \textbf{Tamaño y estratificación muestral objetivo ($N$)}: Se establece una muestra total de $N = 200$ anuncios de empleo curados y estructurados recopilados en Chile (2024–2026). Para capturar adecuadamente la diversidad de la demanda laboral local, la muestra se divide formalmente en dos estratos:
    \begin{itemize}
        \item \textbf{Sector Tecnológico (Mundo TI):} Empresas cuyo giro principal es el desarrollo de software, servicios de consultoría TI, telecomunicaciones y startups tecnológicas.
        \item \textbf{Sector No-TI (Industria Usuaria):} Empresas de sectores tradicionales (banca e instituciones financieras, minería, retail, salud, energía y servicios públicos) que requieren perfiles informáticos para sus departamentos de tecnología interna y transformación digital.
    \end{itemize}
    \item \textbf{Ventana temporal de captura}: Avisos de trabajo publicados activamente entre enero de 2024 y junio de 2026.
    \item \textbf{Categorías de cargos consideradas}: La selección abarca los perfiles profesionales más representativos del mercado de TI en Chile:
    \begin{itemize}
        \item Desarrollador de Software / Ingeniero de Software (\textit{Backend}, \textit{Frontend}, \textit{Fullstack}).
        \item Ingeniero / Arquitecto de Datos (\textit{Data Engineer / Architect}).
        \item Científico de Datos / Analista de Inteligencia Artificial (\textit{Data Scientist / AI Specialist}).
        \item Ingeniero de Infraestructura y Operaciones (\textit{DevOps / Cloud Engineer}).
        \item Ingeniero en Ciberseguridad y Redes.
        \item Arquitecto de Software y Sistemas.
    \end{itemize}
    \item \textbf{Criterios de inclusión y exclusión}:
    \begin{itemize}
        \item \textit{Inclusión}: Ofertas laborales localizadas en Chile (presencial, híbrido o remoto), con descripción detallada de funciones, requisitos técnicos explícitos y nivel de experiencia (junior, semi-senior, senior).
        \item \textit{Exclusión}: Avisos duplicados, ofertas de pasantías o prácticas profesionales no remuneradas, solicitudes con descripciones genéricas sin requerimientos tecnológicos concretos, y vacantes orientadas puramente a ventas o gestión comercial sin exigencia técnica.
    \end{itemize}
    \item \textbf{Protocolo de curaduría y lematización}: Extracción de textos mediante estructuración de secciones de ``Requisitos'' y ``Responsabilidades'', seguida por filtrado de palabras vacías laborales (ej. ``empresa'', ``beneficios'', ``postular'') y extracción de entidades tecnológicas (lenguajes, frameworks, metodologías).
    \item \textbf{Cumplimiento de Términos de Servicio (ToS)}: La captura de datos se realiza en estricto cumplimiento con los Términos de Servicio de LinkedIn, utilizando información de acceso público de forma manual o mediante APIs de consulta de vacantes públicas para fines exclusivamente de investigación académica sin fines de lucro, sin scraping masivo ni almacenamiento de datos personales identificables de reclutadores o postulantes.
\end{enumerate}

% Anexo 5
\subsection{Flujo de procesamiento metodológico}
\label{app:flujo_metodologico}

\begin{table}[H]
    \centering
    \small
    \renewcommand{\baselinestretch}{1.0}\selectfont
    \caption{Detalle de entradas, procesos y salidas del flujo metodológico mínimo.}
    \label{tab:flujo_metodologico_detalle}
    \captionsetup{skip=8pt}
    \begin{tabularx}{\textwidth}{l >{\raggedright\arraybackslash}X >{\raggedright\arraybackslash}X}
        \toprule
        \textbf{Fase} & \textbf{Entrada / Fuentes} & \textbf{Resultado / Salida} \\
        \midrule
        1. Recolección & Mallas, PDF de asignaturas, estándares, LinkedIn & Documentos crudos en base de datos \\
        2. Limpieza & Textos crudos extraídos & Textos planos limpios en minúsculas \\
        3. Normalización & Nombres de asignaturas heterogéneos & Nombres estandarizados y homologados \\
        4. Construcción Corpus & Textos normalizados y filtrados & Base de datos del corpus estructurado \\
        5. Rep. Vectorial & Corpus estructurado & Vectores densos (embeddings multilingües) \\
        6. Reducción Dim. & Vectores de alta dimensionalidad & Vectores en espacio de menor dimensión \\
        7. Clustering & Vectores de asignaturas/perfiles & Clústeres (K-Means, Jerárquico) y Outliers (DBSCAN) \\
        8. Similitud & Vectores y centroides de grupos & Índices de brecha y alineación semántica \\
        9. Visualización & Métricas de alineación y clústeres & Gráficos interactivos y reportes visuales \\
        \bottomrule
    \end{tabularx}
\end{table}

% Anexo 6
\subsection{Planificación temporal del proyecto (Carta Gantt)}
\label{app:gantt}

\begin{table}[H]
    \centering
    \small
    \renewcommand{\baselinestretch}{1.0}\selectfont
    \caption{Carta Gantt con la distribución temporal de las actividades (en semanas).}
    \label{tab:carta_gantt}
    \captionsetup{skip=8pt}
    \begin{tabularx}{\textwidth}{>{\raggedright\arraybackslash}p{6cm} c c c c c c c c}
        \toprule
        \textbf{Actividad / Semana} & \textbf{1-2} & \textbf{3-4} & \textbf{5-6} & \textbf{7-8} & \textbf{9-10} & \textbf{11-12} & \textbf{13-14} & \textbf{15-16} \\
        \midrule
        Fase 1: Recolección y comprensión & \cellcolor{gray!40}X & \cellcolor{gray!40}X & & & & & & \\
        Fase 2: Preparación de datos & & \cellcolor{gray!40}X & \cellcolor{gray!40}X & & & & & \\
        Fase 3: Modelamiento y embeddings & & & \cellcolor{gray!40}X & \cellcolor{gray!40}X & \cellcolor{gray!40}X & & & \\
        Fase 4: Evaluación y métricas & & & & & \cellcolor{gray!40}X & \cellcolor{gray!40}X & & \\
        Fase 5: Despliegue (Dashboard) & & & & & & \cellcolor{gray!40}X & \cellcolor{gray!40}X & \\
        Redacción informe y validación & \cellcolor{gray!20}X & \cellcolor{gray!20}X & \cellcolor{gray!20}X & \cellcolor{gray!20}X & \cellcolor{gray!20}X & \cellcolor{gray!20}X & \cellcolor{gray!40}X & \cellcolor{gray!40}X \\
        \bottomrule
    \end{tabularx}
\end{table}



\end{document}